\begin{figure}[htbp]
  \centering
  \resizebox{\textwidth}{!}{%
  \begin{tikzpicture}[
    >=Latex,
    font=\small,
    node distance=11mm and 12mm,
    class/.style={
      draw,
      rounded corners=3pt,
      align=left,
      fill=gray!8,
      minimum width=38mm
    },
    support/.style={
      draw,
      rounded corners=3pt,
      align=left,
      fill=blue!6,
      minimum width=34mm
    },
    arrow/.style={-{Latex[length=2.3mm]}, thick}
  ]
    \node[class] (settings) {
      \textbf{GenerationSettings}\\
      key\\
      bars\\
      allowed durations\\
      range\\
      form plan\\
      harmony plan\\
      texture\\
      voice profile\\
      chorale plan
    };

    \node[support, left=of settings] (motif) {
      \textbf{Motif}\\
      events: NoteEvent[]
    };

    \node[support, above right=of settings] (form) {
      \textbf{FormPlan}\\
      sections: FormSection[]
    };

    \node[support, right=of settings] (harmony) {
      \textbf{HarmonyPlan}\\
      spans: HarmonySpan[]
    };

    \node[support, right=of harmony] (span) {
      \textbf{HarmonySpan}\\
      start bar / beat\\
      end bar / beat\\
      Roman symbol
    };

    \node[support, above left=of settings] (voice) {
      \textbf{VoiceProfile}\\
      range\\
      tessitura\\
      clef hint
    };

    \node[class, below=24mm of settings] (melody) {
      \textbf{Melody}\\
      events: NoteEvent[]\\
      key\\
      metadata\\
      clef
    };

    \node[class, right=of melody] (chorale_result) {
      \textbf{ChoraleScore}\\
      soprano\\
      alto\\
      tenor\\
      bass\\
      harmony plan
    };

    \node[support, left=of melody] (event) {
      \textbf{NoteEvent}\\
      scale step\\
      duration\\
      chromatic adjustment\\
      rest flag
    };

    \draw[arrow] (motif) -- (settings);
    \draw[arrow] (form) -- (settings);
    \draw[arrow] (harmony) -- (settings);
    \draw[arrow] (span) -- (harmony);
    \draw[arrow] (voice) -- (settings);
    \draw[arrow] (settings) -- node[left] {used in generating} (melody);
    \draw[arrow] (settings) -- node[right] {also drives} (chorale_result);
    \draw[arrow] (event) -- node[above] {used in} (melody);
    \draw[arrow] (motif) -- node[left] {reused as} (event);
  \end{tikzpicture}%
  }
  \caption{Simplified object model of the third iteration. Compared with the earlier iterations, musical information is no longer limited to notes and phrases, but also includes beat-resolved harmony spans, texture settings, voice profiles, rests, and a separate chorale result type.}
  \label{fig:iter3_object_model}
\end{figure}
